\section{Gaussian process 把核变成函数先验}
\label{sec:13-Machine-Learning/09-Kernel-Methods-and-Gaussian-Processes/05}

沿一条路布了三个空气质量监测点，环保部门要你估计中间某个没有设站位置的浓度，据此决定要不要在那里加设一台监测仪。上一节的核岭回归给得出一个数：\(18.4\)。对方追问，这个数误差有多大。你能给的只有交叉验证的平均误差，可那是整条路上的平均，不是这个位置的。更麻烦的是，对方接着又问了一个离所有站点都很远的位置——模型照样输出一个数，语气一样笃定。

点估计的结构里没有``这里我不确定''的位置。这一节换一个视角看同一个核矩阵：把 \(K\) 读成函数值之间的协方差，预测就从一个数变成一个分布，均值与方差一起产出。为此要把先验说清楚，而收益也有明确边界：方差回答的是``这里有没有数据''，不回答``模型有没有选对''。

\subsection{把核读成函数值的联合协方差}

先忘掉数据。假设你对函数 \(f(x)\) 长什么样有一个先验信念：相近的输入上函数值应该相近，远离的输入上可以相差很大。怎样把这句直觉变成数学？选一个正定核 \(k(x,x')\)，让它的值直接充当 \(f(x)\) 与 \(f(x')\) 的协方差。常见的选择是平方指数核（也就是前面的 RBF 核换一种参数化）：

\[
k(x,x')=\sigma_f^2\exp\left(-\frac{(x-x')^2}{2\ell^2}\right).
\]

两个参数各控制一件事。\(\ell\) 是长度尺度，规定多远才算``相近''：\(\ell=1\) 时间隔两三个单位相关性就明显下降，\(\ell=10\) 时函数值在几十个单位内仍高度相关，函数因而更平滑。\(\sigma_f^2\) 是信号方差，规定函数值本身的摆动幅度：\(\sigma_f^2=10\) 意味着你预期 \(f\) 在均值上下三个单位左右起伏。

Gaussian process 的声明只有一句话：对任意有限个输入 \(x_1,\dots,x_n\)，函数值向量服从多元 Gaussian 分布

\[
\bigl(f(x_1),\dots,f(x_n)\bigr)\sim\mathcal{N}(0,K),
\qquad K_{ij}=k(x_i,x_j).
\]

这句话定义的不是一个函数，而是\textbf{函数上的一个概率分布}；它用有限维分布刻画一个无穷维对象，正是\hyperref[sec:06-Random-Processes/01-Foundations/02]{``有限维分布怎样描述一个过程''}讲的那种定义方式，核的半正定性保证这族有限维分布相容。

零均值不是限制，实践中先减去样本均值即可，``回到先验均值''也就成了``回到样本均值''。三个监测点在 \(x=1,3,5\)，取 \(\sigma_f^2=10\)、\(\ell=1.5\)，先验协方差矩阵是

\[
K=
\begin{bmatrix}
10 & 4.11 & 0.29\\
4.11 & 10 & 4.11\\
0.29 & 4.11 & 10
\end{bmatrix}.
\]

相邻两点相距 \(2\)，相关系数 \(0.41\)；相距 \(4\) 时骤降到 \(0.03\)。\(\ell=1.5\) 这句平滑假设就具体地体现在这些数字里。从这个分布可以采样函数：抽一个向量，它的分量就是这些位置上的函数值。\(\ell=1.5\) 的样本平滑而有局部起伏，\(\ell=0.2\) 的样本在很短距离内剧烈震荡。选核就是选先验。

\subsection{观测把先验条件化成后验}

现在加入数据 \((x_1,y_1),\dots,(x_n,y_n)\)，假设每个观测是函数值加上独立测量噪声：

\[
y_i=f(x_i)+\varepsilon_i,\qquad\varepsilon_i\sim\mathcal{N}(0,\sigma_n^2).
\]

\(\sigma_n^2\) 表示你认为测量有多可靠：\(\sigma_n^2=0.01\) 近乎完全信任读数，\(\sigma_n^2=1\) 则只把读数当作大致趋势。三个监测点的读数为 \(y=(10.2,\,14.1,\,18.3)\)，取 \(\sigma_n^2=0.1\)。

关键操作是把观测值与待预测值拼成一个联合 Gaussian 分布。记训练输入为 \(X\)，测试输入为 \(x_*\)，则

\[
\begin{bmatrix}y\\ f_*\end{bmatrix}
\sim\mathcal{N}\left(
0,\
\begin{bmatrix}
K_{XX}+\sigma_n^2 I & K_{X*}\\
K_{*X} & K_{**}
\end{bmatrix}
\right).
\]

剩下的全是\hyperref[sec:05-Probability/06-Common-Distributions/05]{``Gaussian 与多元 Gaussian''}里的条件分布公式：给定 \(y\) 后 \(f_*\) 仍是 Gaussian，

\[
f_*\given y\ \sim\ \mathcal{N}(\mu_*,\ \sigma_*^2),
\qquad
\left\{
\begin{aligned}
\mu_*&=K_{*X}(K_{XX}+\sigma_n^2I)^{-1}y,\\
\sigma_*^2&=K_{**}-K_{*X}(K_{XX}+\sigma_n^2I)^{-1}K_{X*}.
\end{aligned}
\right.
\]

不必被逆矩阵吓到，第一式是对观测值做加权平均，权重由``新点与各训练点的相关性''决定；第二式则是从先验方差 \(K_{**}\) 里扣掉被观测解释掉的那部分。同一套代数在\hyperref[sec:13-Machine-Learning/07-Graphical-and-Sequential-Models/03]{``Kalman filter 是动态系统上的 Gaussian 条件化''}中沿时间反复施行，这里只施行一次。

代入数字。对 \(x_*=4\)，核向量 \(K_{*X}\approx(1.35,\ 8.01,\ 8.01)\)：新点离 \(x_2=3\) 与 \(x_3=5\) 很近，离 \(x_1=1\) 较远。算得

\[
\mu_*\approx 16.90,\qquad \sigma_*^2\approx 0.82,
\]

预测标准差约 \(0.91\)。换成 \(x_*=20\)，三个核值都在 \(10^{-15}\) 量级以下，长度尺度 \(\ell=1.5\) 认为这里与训练数据几乎无关，于是

\[
\mu_*\approx 14.20,\qquad \sigma_*^2\approx 10.00,
\]

均值退回先验均值（即样本均值），方差退回先验方差，与没有观测时一模一样。这正是点估计给不出的那个信号：此处无据可依，别把这个数当结论。

把 \(x_*\) 扫过整条路，就得到下面这张图。

\begin{center}
\begin{tikzpicture}[x=0.95cm,y=0.40cm,>=stealth,font=\small]

\fill[black!10]
  plot[smooth] coordinates {(-1.00,18.23) (-0.50,16.63) (0.00,14.54)
    (0.50,12.30) (1.00,10.87) (1.50,12.00) (2.00,13.23) (2.50,13.97)
    (3.00,14.73) (3.50,16.91) (4.00,18.72) (4.50,19.26) (5.00,18.89)
    (5.50,20.02) (6.00,21.12) (6.50,21.62) (7.00,21.62) (7.50,21.35)
    (8.00,21.03) (8.50,20.79) (9.00,20.64)}
  -- plot[smooth] coordinates {(9.00,8.00) (8.50,8.17) (8.00,8.52)
    (7.50,9.17) (7.00,10.26) (6.50,11.90) (6.00,14.01) (5.50,16.24)
    (5.00,17.63) (4.50,16.42) (4.00,15.09) (3.50,14.27) (3.00,13.47)
    (2.50,11.32) (2.00,9.60) (1.50,9.15) (1.00,9.61) (0.50,8.52)
    (0.00,7.43) (-0.50,6.91) (-1.00,6.87)}
  -- cycle;

\draw[dashed,black!50] (-1.3,14.2) -- (9.5,14.2)
  node[right,font=\footnotesize] {先验均值};

\draw[very thick]
  plot[smooth] coordinates {(-1.00,12.55) (-0.50,11.77) (0.00,10.99)
    (0.50,10.41) (1.00,10.24) (1.50,10.57) (2.00,11.41) (2.50,12.64)
    (3.00,14.10) (3.50,15.59) (4.00,16.90) (4.50,17.84) (5.00,18.26)
    (5.50,18.13) (6.00,17.56) (6.50,16.76) (7.00,15.94) (7.50,15.26)
    (8.00,14.78) (8.50,14.48) (9.00,14.32)};

\fill (1,10.2) circle (2.4pt);
\fill (3,14.1) circle (2.4pt);
\fill (5,18.3) circle (2.4pt);

\draw[->] (-1.3,6.3) -- (9.7,6.3) node[right] {\(x\)};
\draw[->] (-1.3,6.3) -- (-1.3,22.6) node[above] {\(y\)};
\foreach \t in {1,3,5,7}
  \draw (\t,6.3) -- (\t,5.9) node[below,font=\footnotesize] {\t};
\foreach \v in {10,15,20}
  \draw (-1.3,\v) -- (-1.5,\v) node[left,font=\footnotesize] {\v};

\draw[dotted,black!60] (8,7.4) -- (8,21.6);
\node[font=\footnotesize,align=left,anchor=north west] at (5.6,13.2)
  {远离数据：\\带宽回到先验};

\end{tikzpicture}
\end{center}

\emph{阴影是约 \(95\%\) 的预测区间；它在观测点处收紧，在没有数据的地方张开到先验宽度。}

\subsection{后验均值就是核岭回归，解释却不同}

看一眼 \(\mu_*=K_{*X}(K_{XX}+\sigma_n^2I)^{-1}y\)，再看上一节核岭回归的预测

\[
\widehat f(x_*)=K_{*X}(K_{XX}+\lambda nI)^{-1}y.
\]

两式逐字相同，对应关系是 \(\sigma_n^2\leftrightarrow \lambda n\)。同一段代码可以同时是两个模型的实现。差别在解释：核岭回归通过最小化正则化经验风险得到一个点估计，正则参数是控制 RKHS 范数的预算；Gaussian process 从函数先验与噪声模型出发做条件化，得到完整的预测分布，那个参数是观测噪声的方差。要点预测，两者给出同一个数；要``这个预测有多可靠''，只有后者答得上来——因为只有后者的模型里写下了方差。

这也划出一条使用纪律：把核岭回归的 \(\lambda\) 随口称作噪声方差，只有在明确接受了 Gaussian process 那整套概率模型之后才合法。同样的对应在\hyperref[sec:13-Machine-Learning/02-Linear-Models/04]{``Bayesian 线性回归把系数不确定性传到预测''}中出现过：Ridge 的解等于 Gaussian 先验下的后验众数，而后验还多给了一个协方差。代数相等，认知结构不同。

\subsection{核参数可以由边缘似然选出来}

上面直接设定了 \(\ell=1.5\)、\(\sigma_f^2=10\)、\(\sigma_n^2=0.1\)。实践中这些参数通常由数据决定，Gaussian process 的自然做法是最大化边缘似然，即观测 \(y\) 在当前核参数下的对数概率：

\[
\log p(y\given X,\theta)
=-\frac12 y^{\trans}K_y^{-1}y
-\frac12\log\det K_y
-\frac n2\log(2\pi),
\]

其中 \(K_y=K_{XX}+\sigma_n^2I\)，\(\theta\) 收集 \((\ell,\sigma_f^2,\sigma_n^2)\)。

前两项分工明确。\(-\tfrac12 y^{\trans}K_y^{-1}y\) 是数据拟合项：若当前核认为两点高度相关而观测值差异很大，协方差矩阵解释不了数据，这一项的绝对值就变大，模型和数据打架就会在这里显出来。\(-\tfrac12\log\det K_y\) 是复杂度惩罚项：\(\ell\) 很小时所有点被认为几乎不相关，\(K_y\) 接近对角阵，\(\det K_y\) 很大，这一项把过度灵活的先验压下去。两项合力，一项驱使模型贴合数据，一项阻止它用极端灵活性去贴合。与交叉验证相比，边缘似然不需要划分数据，在同一批观测上同时完成拟合与复杂度控制。

它也不是免费的。边缘似然关于 \(\theta\) 通常非凸，可能有多个局部最优——长度尺度很大配大噪声（``趋势加噪音''）与长度尺度很小配小噪声（``什么都是信号''）常常是两个互相竞争的解释，优化落到哪一个取决于初值。最大化边缘似然是在给定核族内部选参数，它不会告诉你核族本身选错了。

\subsection{Cholesky 决定数值稳定，\texorpdfstring{\(O(n^3)\)}{O(n3)} 决定规模}

公式里的 \((K_{XX}+\sigma_n^2I)^{-1}\) 在代码中不应显式求逆，而应走\hyperref[sec:03-Linear-Algebra/06-Symmetric-and-Positive-Definite-Matrices/03]{``Cholesky 分解：把正能量写成平方和''}那条路：

\[
K_y=LL^{\trans},
\]

再解三角方程组 \(Lz=y\) 与 \(L^{\trans}\alpha=z\) 得到 \(\alpha=K_y^{-1}y\)，预测均值即 \(K_{*X}\alpha\)。同一次分解还顺带给出

\[
\log\det K_y=2\sum_i\log L_{ii},
\]

于是预测均值、预测方差与边缘似然共用一次 \(O(n^3)\) 的分解，且比显式求逆数值稳定。

这个 \(O(n^3)\) 与 \(O(n^2)\) 存储与上一节核岭回归的三笔账是同一笔账，因为要分解的本来就是同一个矩阵。\(n\) 超过几千，稠密 Cholesky 就开始吃力；出路也相同：诱导点方法（用 \(m\ll n\) 个虚拟位置概括训练集，把复杂度压到 \(O(nm^2)\)）、随机特征、结构化核。区别在于 Gaussian process 的近似必须同时保住方差的正确性，只把均值近似好是不够的——一个把方差系统性压小的近似，恰好毁掉了使用它的理由。

另有一件工程细节值得记住。核矩阵会因重复输入、极端长度尺度或舍入误差而接近奇异，此时在对角线上加一个微小的 jitter（如 \(10^{-8}\)）是常见保护。Jitter 是数值稳定装置，与统计噪声参数 \(\sigma_n^2\) 不是一回事，把它调大到能明显改变预测，说明模型设定已经出了问题。

\subsection{方差只回答``这里有没有数据''}

现在必须把图中那条阴影带的含义收紧。后验方差

\[
\sigma_*^2=K_{**}-K_{*X}(K_{XX}+\sigma_n^2I)^{-1}K_{X*}
\]

有一个容易被忽略的特点：它\textbf{不含 \(y\)}。观测值是多少、模型拟合得好不好，对方差没有任何影响；只有输入位置 \(x_*\) 与训练输入 \(X\) 的相对位置进得去。换句话说，\(\sigma_*^2\) 度量的是``在当前核所定义的几何下，这个位置离数据有多远''，而不是``当前模型在这个位置错得有多离谱''。图里带宽在观测点处收紧、在两端张开，全部信息来自横轴上点的疏密。

这个性质带来一半好消息。远离数据处方差大，是对的，也正是环保部门需要的信号；如果换一种布点方案能让整条路上的方差都小一些，那么无须采集任何读数就能比较两种方案的优劣——主动学习与实验设计正是这么做的。

坏消息在另一半。方差不反映模型设定是否正确，下面这些误差来源全部落在它的视野之外。核族选错了：你用平方指数核去拟合一条带周期成分的曲线，在两个观测点之间方差照样很小，而均值会平滑地穿过本该出现的波峰。核参数估偏了：边缘似然给出的 \(\ell\) 只是估计，若它比真实相关尺度大一倍，模型会误以为远处的点也被观测覆盖，于是在数据稀疏处报出过小的方差。观测噪声模型错了：假设 Gaussian 噪声而真实噪声重尾时，个别离群读数会把均值拽偏，方差却不会相应放大。部署分布变了：新样本来自训练期从未出现的工况，只要它在输入空间里恰好落在训练点附近，核就认定它们相似，方差照样很小——这与\hyperref[sec:13-Machine-Learning/01-Learning-Problems-and-Evaluation/04]{``数据划分、泄漏与分布偏移''}中提醒的失效模式是同一类。

于是有一条不对称的读法。方差大的地方，预测不可靠，这半句永远成立，因为没有数据就是没有数据。方差小的地方预测可靠，这半句只在核族、核参数与数据分布都正确的前提下成立。后验方差是``给定这个模型''的诚实报告，不是现实世界总不确定性的保险单。要检验后者，得把预测区间拿到真实数据上校准，看约 \(95\%\) 的区间是否真的覆盖了约 \(95\%\) 的结果——\hyperref[sec:11-Mathematical-Statistics/07-Statistical-Decision-and-Reliable-Prediction/04]{``Proper scoring rule 与校准衡量概率是否诚实''}提供了做这件事的工具，而\hyperref[sec:11-Mathematical-Statistics/07-Statistical-Decision-and-Reliable-Prediction/05]{``Conformal prediction 用可交换性换有限样本覆盖''}给出一条不依赖核设定正确的补救路线。

\subsection{选核是建模，不是调参}

最后回到核本身。平方指数核无限可微，产生的样本函数极其光滑，真实函数若有突变或拐角，它就不合适。Matérn 核族多出一个光滑度参数 \(\nu\)：\(\nu=1/2\) 对应连续但处处不可微的路径，\(\nu=3/2\) 与 \(\nu=5/2\) 分别对应一次与两次可微，\(\nu\to\infty\) 回到平方指数核。周期核捕捉季节性重复。核还可以相加相乘——平方指数核加周期核，声明的是``既有长期趋势又有周期波动''。

这些旋钮改的是同一件事：什么样的函数算平滑。这句声明在本单元里出现过三次，形态各不相同。第二节里它是特征空间的形状，第三节里它是 RKHS 范数对函数的收费标准，这一节里它是函数上的一个概率分布。三种说法指向同一个 \(k\)，而选择 \(k\) 始终是建模决定，不是超参数搜索能代劳的事——搜索只能在你圈定的核族里挑，圈错了族，边缘似然与验证误差都会在错误的范围内给出一个漂亮的最优值。

这也是\hyperref[ch:117]{``核方法与 Gaussian process：把函数当作随机对象''}这一单元最后留下的东西：从``数据只以内积出现''这句观察出发，一路推到函数上的分布，途中每一步都是定理；唯独第一步——写下哪个 \(k\)——始终由人负责，而后面所有的均值、方差与置信区间，都是在这个人为选择之内说话。
